
%
% Jeu de couleurs de l'établissement d'exemple Ecole1
% (mêmes commandes que les autres jeux : \TR, \TV, \TB… ; voir style/latex/couleurs_roles.tex)
%

% Couleurs principales
% Blanc
\newcommand {\TW}[1] {\textcolor{white}{#1}}
% Encre : 29 29 29
\definecolor{noirC}{RGB}{29,29,29}
\newcommand{\TN}[1]{\textcolor{noirC}{#1}}
% Rouge : 223 35 29
\definecolor{rougeF}{RGB}{223,35,29}
\newcommand{\TR}[1]{\textcolor{rougeF}{#1}}
% Violet : 71 10 104
\definecolor{violetR}{RGB}{71,10,104}
\newcommand{\TV}[1]{\textcolor{violetR}{#1}}
% Orange vif : 252 76 2
\definecolor{orangeB}{RGB}{252,76,2}
\newcommand{\TOB}[1]{\textcolor{orangeB}{#1}}
% Jaune : 255 205 0
\definecolor{jauneFA}{RGB}{255,205,0}
\newcommand{\TJ}[1]{\textcolor{jauneFA}{#1}}
% Cyan : 100 204 201
\definecolor{cyanI}{RGB}{100,204,201}
\newcommand{\TC}[1]{\textcolor{cyanI}{#1}}
% Bleu : 77 98 201
\definecolor{bleuE}{RGB}{77,98,201}
\newcommand{\TB}[1]{\textcolor{bleuE}{#1}}
% Rouge brun : 124 38 41
\definecolor{rougeS}{RGB}{124,38,41}
\newcommand{\TRS}[1]{\textcolor{rougeS}{#1}}
% Mauve : 159 40 66
\definecolor{mauveSI}{RGB}{159,40,66}
\newcommand{\TM}[1]{\textcolor{mauveSI}{#1}}
% Orange : 255 143 28
\definecolor{orangeVE}{RGB}{255,143,28}
\newcommand{\TO}[1]{\textcolor{orangeVE}{#1}}
% Gris secondaire : 221 221 221
\definecolor{grisS}{RGB}{221,221,221}
\newcommand{\TG}[1]{\textcolor{grisS}{#1}}

% Couleurs d'accent
% Orange
\definecolor{orangeCent}{RGB}{255,97,27}
\newcommand{\TOC}[1]{\textcolor{orangeCent}{#1}}
% Rose
\definecolor{roseCent}{RGB}{247,66,159}
\newcommand{\TRC}[1]{\textcolor{roseCent}{#1}}
% Violet
\definecolor{violetCent}{RGB}{122,55,228}
\newcommand{\TVC}[1]{\textcolor{violetCent}{#1}}



%PERSO
\newcommand {\tw}[1] {\textcolor{white}{#1}}

%%%%%%%%%%%%%%%%%%%%%%%%%%%%%%%%%%%%%%%%%%%%%%%%%%%%%%%%%%%%%%%%%%%%%%%%%%%%%%
%% HARMONISATION (proposition v2) — à partir d'ici, mêmes noms dans tous les jeux
%%%%%%%%%%%%%%%%%%%%%%%%%%%%%%%%%%%%%%%%%%%%%%%%%%%%%%%%%%%%%%%%%%%%%%%%%%%%%%
% Les commandes ci-dessus restent valables (anciennes diapos). Attention :
% \TV = violet ici, mais vert dans New_BD_ColorScheme : préférer \TVi / \TVe.

% Teintes communes (couleur la plus proche de la charte)
% Gris foncé : absent de la charte
\definecolor{grisF}{RGB}{95,99,104}
\newcommand{\TS}[1]{\textcolor{grisF}{#1}}    % gris foncé
\newcommand{\TVi}[1]{\textcolor{violetR}{#1}} % violet
% Vert : absent de la charte, remplacé par le cyan
\newcommand{\TVe}[1]{\textcolor{cyanI}{#1}}   % vert

% Couleurs de rôle (commandes \TDef, \TPro, \TRes, \TAtt, \TMC, \TExg dans couleurs_roles.tex)
\colorlet{roleDef}{bleuE}     % rappel, définition
\colorlet{rolePro}{violetR}   % proposition, modèle
\colorlet{roleRes}{roseCent}   % résultat à retenir
\colorlet{roleAtt}{rougeF}    % attention, erreur
\colorlet{roleExg}{violetR}   % mise en exergue (encadré)
\colorlet{roleMC}{violetR}    % couleur principale : bandeau, titres
\input{couleurs_roles}
